\documentclass[conference]{IEEEtran}

\usepackage[T1]{fontenc}
\usepackage[utf8]{inputenc}
\usepackage{amsmath}
\usepackage{booktabs}
\usepackage{cite}
\usepackage{graphicx}
\usepackage{hyperref}
\usepackage{url}

\setlength{\textfloatsep}{6pt plus 1pt minus 2pt}
\setlength{\dbltextfloatsep}{6pt plus 1pt minus 2pt}
\setlength{\floatsep}{4pt plus 1pt minus 1pt}
\setlength{\dblfloatsep}{4pt plus 1pt minus 1pt}
\setlength{\abovecaptionskip}{2pt}
\setlength{\belowcaptionskip}{0pt}

\hypersetup{
  colorlinks=true,
  linkcolor=blue,
  citecolor=blue,
  urlcolor=blue
}

\title{MICO: Compiler-Gated Structured Module Composition for LLM-Assisted RTL Integration}

\author{
\IEEEauthorblockN{Anonymous Author(s)}
\IEEEauthorblockA{Paper under double-blind review}
}

\begin{document}

\maketitle

\begin{abstract}
LLMs can generate small Verilog fragments, but RTL integration remains brittle because correctness depends on port binding, direction, width, handshake protocol, clock/reset domains, CDC/RDC handling, adapters, and wrappers. MICO reframes LLM-assisted integration as compiler-gated structured module composition: models propose typed interface graphs and bounded repair patches, while a Rust compiler validates the graph and emits traceable SystemVerilog, SVA skeletons, JSON IR, and diagnostics. The artifact passes an expanded deterministic benchmark boundary with 83/83 public-development ModuleComposeBench tasks, 40/40 held-out tasks, and a 30/30 realism supplement, including 46/46 expanded deterministic public positive simulations, 40/40 expanded deterministic single-clock public bounded formal smoke checks, 11/11 expanded deterministic structural/generic public Yosys QoR rows, and a Vivado out-of-context subset covering 21 reference-enabled split rows through 12 unique task pairs. In an authenticated v3 LLM matrix bound to the locked pre-expansion 62-task public and 20-task held-out manifests, schema-guided MICO JSON AST repair reaches 35--36/36 public positive compiler/lint passes, 9--10/10 held-out positive passes, 25--26/26 public unsafe rejections, and 10/10 held-out unsafe rejections, while direct Verilog and SystemVerilog-interface prompts lint at most 5/36 public positives and 2/10 held-out positives. This is a bounded tested-profile claim, not broad autonomous repair: recorded repair wins come from an explicitly recorded deterministic compiler-feedback fallback for one adapter-instance error pattern, and MICO does not claim exhaustive formal coverage, CDC correctness proof, routed timing closure, arbitrary LTL, or generalization to untested models.
\end{abstract}

\begin{IEEEkeywords}
RTL integration, large language models, hardware description languages, static checking, module composition, SystemVerilog
\end{IEEEkeywords}

\section{Introduction}

Much top-level RTL and subsystem code does not implement new datapath logic; it connects existing modules, wrappers, channels, buses, and clock/reset infrastructure. A ready/valid edge in a diagram expands into payload, valid, ready, sideband, reset, and sometimes CDC or protocol-adapter wiring. For LLMs this is a poor target: free-form RTL requires long-range consistency over names, directions, widths, protocols, domains, resets, and validation collateral.

The resulting failures are systematic rather than cosmetic. In the benchmark
records, unstructured prompts most often fail by omitting ports, reversing
producer/consumer direction, crossing clock domains without an explicit
adapter, emitting syntactically plausible but semantically rejected wrappers, or
returning JSON that cannot be machine-checked. MICO's causal hypothesis is that
the target representation should expose these obligations as typed graph edges
and contract checks before any RTL is accepted.

MICO takes a narrower position than general RTL generation benchmarks~\cite{verilogeval2023,rtlcoder2023,openllmrtl2025}, domain-adapted chip LLMs~\cite{chipnemo2023}, and recent checked prompt languages~\cite{cppl2026}. The LLM proposes a module-composition graph and optional repair patch; the compiler is the authority. The compiler checks interfaces, roles, widths, protocols, clock domains, adapters, and v0 ready/valid contracts before emitting RTL and traceability artifacts.

This paper makes five contributions:

\begin{enumerate}
  \item a typed module-interface-contract representation that separates LLM proposal from compiler authority;
  \item the MICO language and Rust compiler path for checked graphs, diagnostics, JSON AST repair, SV/SVA emit, and traceability;
  \item ModuleComposeBench, with public-development, held-out, and supplemental realism manifests, positive/negative tasks, case studies, simulation, formal smoke, and QoR collateral;
  \item an authenticated full-matrix evaluation showing that schema-guided JSON AST composition and a bounded compiler-feedback fallback improve tested LLM integration correctness and unsafe rejection over direct RTL, SystemVerilog-interface, and MICO-source prompting;
  \item a reproducible artifact with Docker-first validation, schema gates, release hashes, a Vivado host-tool exception, and reviewer bundle metadata.
\end{enumerate}

The positive LLM claim is deliberately narrow. It covers the tested profiles,
prompts, manifests, schemas, and compiler gates; it does not make the LLM a
source of correctness. The artifact does not replace HDLs or downstream
verification. It attacks the integration layer where engineers already rely on
structure, naming discipline, and review checklists, and where LLMs are most
exposed to textual consistency failures.

\section{Motivation: Connection Entropy}

We use connection entropy to describe the amplification from semantic design intent to primitive RTL connections:

\begin{equation}
  CE = \frac{\text{primitive signal edges}}{\text{semantic interface edges}}.
\end{equation}

A designer may intend a simple chain from a producer through a FIFO to a consumer. In primitive RTL, that chain expands into many individual signal connections. As connection entropy grows, LLMs are more likely to omit ports, reverse ready/valid direction, mismatch payload widths, confuse reset polarity, cross clock domains without synchronization, generate inconsistent parameters, or produce wrappers that simulate only because the testbench misses the actual protocol bug.

Prompt engineering can reduce some failures, but it does not change the underlying problem: free-form RTL text generation leaves too many coupled correctness obligations in one target string. MICO instead constrains the output space. The model proposes a structured design graph, adapter plan, and contract skeleton. The compiler rejects unsafe connections and returns structured diagnostics that can drive a repair prompt.

The design principles are direct:

\begin{itemize}
  \item LLMs produce candidates, not trusted RTL.
  \item Direction, width, protocol, and domain mismatches are compile-time errors.
  \item CDC and RDC crossings require explicit adapters.
  \item Automatically inserted buffers, pipelines, width adapters, or CDC FIFOs must appear in an elaboration report.
  \item Verification collateral is generated from interface contracts.
\end{itemize}

This framing distinguishes two kinds of complexity. \emph{Local RTL complexity} is the logic inside a module: arithmetic, state machines, memories, and datapaths. \emph{Composition complexity} is the set of obligations required to connect already-existing modules safely. MICO focuses on the latter. A direct RTL prompt must solve both at once. A MICO prompt only proposes a semantic graph; the compiler expands that graph into primitive wiring and rejects illegal edges before RTL tools are asked to elaborate it.

\begin{table}[t]
  \centering
  \caption{Typical composition failures targeted by MICO}
  \begin{tabular}{lll}
    \toprule
    Failure & Free-form RTL symptom & MICO gate \\
    \midrule
    Direction & ready/valid reversed & role check \\
    Width & silent truncation & field width check \\
    Protocol & ad hoc handshake & interface contract \\
    CDC & direct cross-domain wire & adapter check \\
    Reset & mixed polarity/domain & clockdom metadata \\
    \bottomrule
  \end{tabular}
\end{table}

\section{Related Work}

\subsection{LLMs for RTL}

VerilogEval, RTLCoder, OpenLLM-RTL, and ChipNeMo show that LLMs can be evaluated and adapted for hardware tasks~\cite{verilogeval2023,rtlcoder2023,openllmrtl2025,chipnemo2023}. These works motivate better benchmarks and domain context, but the artifact emitted by a model is still often RTL text. MICO is complementary: it does not assume the model has become trustworthy. It narrows the model's output to a checked composition representation and uses compiler diagnostics plus RTL tools as the acceptance boundary.

\subsection{Interface Abstractions}

SystemVerilog interfaces, Chisel bundles and bulk connections, Amaranth wiring, and CIRCT ESI all demonstrate that interface abstraction has practical value~\cite{chisel_interfaces,amaranth_wiring,circt_esi}. MICO differs in its objective. It targets the integration of existing RTL/IP modules, not complete hardware description. It restricts LLM output to structured AST/JSON forms, not arbitrary host-language code. It also treats clock/reset domains, protocol contracts, adapter synthesis, and formal collateral as part of the interface type system.

CPPL is the closest recent work~\cite{cppl2026}. If MICO merely claimed ``structured language plus CIRCT,'' it would overlap too much with CPPL. The boundary is narrower: CPPL targets LLM-friendly circuit generation, while MICO targets LLM-friendly module interconnect and integration. Its central objects are interfaces, clock domains, contracts, adapters, and compose graphs.

CIRCT provides a natural downstream target. The HW dialect aligns with Verilog modules, ESI provides typed channels and valid/ready wrapping, and Verif/LTL can represent assertions and temporal properties~\cite{circt_hw,circt_esi}. Anvil and AutoSVA further motivate treating timing safety and module interaction properties as language or tooling concerns rather than informal comments~\cite{anvil2025,autosva2021}.

\subsection{Verification and Contracts}

MICO's contract layer is not intended to replace simulation, formal verification, lint, CDC analysis, or synthesis. It is a front-end discipline for making obligations explicit early. This is aligned with AutoSVA's focus on module interaction properties and Anvil's emphasis on timing-safety as a language concern~\cite{autosva2021,anvil2025}. The difference is again the LLM boundary: contracts must be available as structured compiler artifacts so a repair loop can preserve or refine them instead of deleting them during a free-form rewrite.

\section{MICO Language Core}

MICO v0 has five core objects: clock domains, interfaces, extern modules, adapters, and compose graphs. A connection is legal only if both instances and ports exist, the source role is an output, the destination role is an input, the interface identities match or an explicit adapter exists, the clock domains match or an explicit CDC adapter exists, and the attached contracts satisfy refinement obligations.

\begin{verbatim}
clockdom Sys(clk, rst);

interface StreamU32 @Sys {
  producer payload:u32, valid:bool;
  consumer ready:bool;
  contract stable_payload:
    valid -> stable(payload) until ready;
}

compose Top @Sys {
  inst p: Producer;
  inst f: Fifo;
  inst c: Consumer;
  connect p.tx -> f.input;
  connect f.output -> c.rx;
}
\end{verbatim}

If any condition fails, the compiler must reject the graph and produce a structured diagnostic. Such diagnostics are designed to be fed into an LLM repair loop as JSON, making repairs local and auditable rather than free-form rewrites of generated RTL.

\subsection{Static Semantics}

The v0 checker enforces a conservative subset of the intended type system. Let an endpoint be a tuple
\[
  e = (\text{instance}, \text{port}, \text{interface}, \text{role}, \text{domain}, \text{fields}, \text{contracts}).
\]
A direct connection \(e_s \rightarrow e_t\) is legal only if \(e_s\) is a producer, \(e_t\) is a consumer, the interface protocol agrees, the clock/reset domain agrees, and each required field has a compatible width and direction. If any of these predicates fail, the graph is rejected unless an explicit adapter node states how the mismatch is resolved.

\begin{table}[t]
  \centering
  \caption{Representative MICO v0 diagnostics}
  \begin{tabular}{lll}
    \toprule
    Check & Example & Repair action \\
    \midrule
    name & duplicate instance & rename or delete \\
    endpoint & unknown port & bind declared port \\
    direction & consumer to consumer & reverse edge \\
    width & u32 to u64 & insert width adapter \\
    domain & Sys to Async & insert CDC adapter \\
    contract & missing obligation & attach contract \\
    \bottomrule
  \end{tabular}
\end{table}

This design intentionally avoids silent convenience. MICO does not insert a CDC FIFO or truncating width adapter merely because doing so would make the wrapper compile. A model or user must declare the adapter explicitly, and the compiler records it in the typed IR and generated report.

\section{Rust Compiler Framework}

MICO is implemented as a Rust workspace with four crates:

\begin{itemize}
  \item \texttt{mico\_ir}: AST, typed IR, diagnostics, and semantic checks.
  \item \texttt{mico\_frontend}: lexer, parser, source maps, and recovery.
  \item \texttt{mico\_codegen}: JSON IR, SystemVerilog, SVA skeleton, and traceability emission.
  \item \texttt{mico\_cli}: \texttt{parse}, \texttt{check}, \texttt{build}, \texttt{dump-ir}, \texttt{emit-sv}, \texttt{emit-sva}, \texttt{emit-trace}, \texttt{verify}, and \texttt{report}.
\end{itemize}

The compiler pipeline is deterministic:

\begin{enumerate}
  \item The frontend lexes and parses MICO source into an AST with source locations.
  \item The IR layer resolves declarations, instances, endpoints, domains, adapter kinds, and contract metadata.
  \item The checker rejects duplicate declarations, unknown endpoints, illegal directions, unsupported adapter kinds, width mismatches, domain mismatches, and missing contract obligations.
  \item Code generation emits versioned JSON IR, SystemVerilog wrappers, SVA skeletons, and traceability reports from checked IR only.
  \item The CLI can emit human-readable or JSON diagnostics so that downstream prompt repair does not depend on terminal text scraping.
\end{enumerate}

Figure~\ref{fig:mico-pipeline} shows the acceptance boundary used by both the
CLI and benchmark harness.

\begin{figure*}[t]
  \centering
  \fbox{\begin{minipage}{0.96\textwidth}
    \centering
    LLM/user proposal
    \(\rightarrow\) MICO source or JSON AST
    \(\rightarrow\) schema/parser
    \(\rightarrow\) typed IR and static checks
    \(\rightarrow\) SV/SVA/traceability JSON
    \(\rightarrow\) Docker EDA and benchmark gates
    \par\smallskip
    Compiler diagnostics
    \(\rightarrow\) bounded JSON repair patch or deterministic fallback
    \(\rightarrow\) re-check before acceptance
  \end{minipage}}
  \caption{MICO compiler-gated pipeline. The model proposes structured data; schemas, the Rust checker, and RTL/EDA gates decide acceptance.}
  \label{fig:mico-pipeline}
\end{figure*}

The first implementation stage intentionally avoids a full MLIR binding. It emits conservative SystemVerilog wrappers, structured JSON IR, SVA skeletons, and traceability reports. A later stage can lower typed IR into CIRCT HW modules, ESI channels, and Verif/LTL assertions. Rust's enums, pattern matching, ownership model, and test ecosystem fit compiler passes and diagnostics.

\section{LLM Collaboration Workflow}

The MICO collaboration loop has five stages:

\begin{enumerate}
  \item inventory modules, ports, protocols, domains, and adapters;
  \item ask the model for a structured graph or JSON AST;
  \item run compiler checks and return structured diagnostics;
  \item ask for a minimal JSON patch when repair is enabled;
  \item run RTL validation and emit traceability from checked artifacts.
\end{enumerate}

This separates proposal from validation. The LLM never becomes the source of correctness; its output is constrained to data structures that the compiler can reject, repair, and compare across model profiles.

The provider policy separates cost from correctness. Smoke tests and prompt development use low-cost profiles first; higher-cost profiles are reserved for candidates that pass structured-output, compiler, and RTL gates.

The provider harness records validate-only and authenticated runs as redacted JSON artifacts: profile, model identifier, prompt hash, request settings, JSON validity, token/cost metadata when available, repair turns, compiler attachment, and EDA attachment. It records whether an API key was present but never stores or prints the key. The batch runner applies this to Direct Verilog, SystemVerilog-interface, MICO source, MICO JSON AST, and JSON AST repair baselines over the benchmark manifest.

\section{Adapter Synthesis}

Adapters are explicit first-class artifacts. Width adapters, protocol bridges, pipeline stages, buffers, and CDC FIFOs must appear in the typed IR and elaboration report. A clock-domain crossing is illegal by default unless the graph declares a verified CDC adapter. Similarly, protocol adaptation is illegal unless an adapter in the library states the source and destination contracts it preserves.

This design is especially important for LLM-assisted flows. If the model silently inserts glue logic, the result may compile while changing latency, backpressure, ordering, or reset semantics. MICO forces these choices into named objects that can be audited, tested, and benchmarked.

\section{Evaluation Plan: ModuleComposeBench}

ModuleComposeBench evaluates whether a model can integrate existing RTL/IP modules, not merely generate standalone Verilog. Tasks include wrapper generation, stream chaining, bus attachment, width adaptation, protocol bridging, CDC insertion, and unsafe ambiguity rejection.

Metrics include first-pass compose success, repair turns, lint pass, simulation pass, formal pass, unsafe rejection rate, QoR delta, and connection entropy reduction. Baselines include direct Verilog prompting, SystemVerilog-interface prompting, MICO AST prompting without compiler-feedback repair, and MICO AST prompting with compiler-feedback repair.

\begin{table}[t]
  \centering
  \caption{Evaluation dimensions}
  \begin{tabular}{lll}
    \toprule
    Dimension & Direct RTL & MICO \\
    \midrule
    Output & RTL text & Typed compose graph \\
    Error boundary & Simulation/review & Compiler + RTL tools \\
    Repair unit & Code rewrite & JSON AST patch \\
    CDC handling & Often implicit & Explicit adapter \\
    \bottomrule
  \end{tabular}
\end{table}

\subsection{Benchmark Protocol}

Each benchmark task contains a module inventory, protocol/interface library, optional adapter library, natural-language integration request, expected safety constraints, and reference validation collateral. A run records the model profile, prompt hash, first-pass compiler result, repair turns, lint result, simulation/formal result when available, unsafe rejection result for negative tasks, and generated artifacts. The same task is evaluated across baselines so that improvements can be attributed to the representation and repair loop rather than task selection.

\begin{table}[t]
  \centering
  \caption{ModuleComposeBench task strata}
  \begin{tabular}{lll}
    \toprule
    Level & Focus & Example \\
    \midrule
    L1 & direct streams & FIFO chain \\
    L2 & widths/params & width adapter \\
    L3 & latency/protocol seeds & skid/pipeline checks \\
    L4 & CDC/RDC & async FIFO \\
    L5 & bus/register seeds & control/status wrapper \\
    L6 & subsystem seeds & multi-block chain \\
    \bottomrule
  \end{tabular}
\end{table}

\subsection{Current Seed Results}

The current artifact includes 57 deterministic seed tasks: 31 positive composition tasks and 26 negative rejection tasks. The manifest spans L1--L6, including same-domain stream wiring, width adaptation, latency/backpressure seeds, CDC/RDC adapter tasks, bus/register wrapper seeds, subsystem seeds, and unsafe examples for direct CDC, width mismatches, reversed direction, missing or invalid contract guarantees, unknown declarations, duplicate declarations, protocol mismatches, and adapter misuse. The L3/L5/L6 entries are still seed approximations over the committed smoke RTL collateral rather than dedicated subsystem case studies, but they validate that the compiler, JSON AST path, ready/valid v0 contract checks, SystemVerilog/SVA emitters, traceability report, Verilator lint, Icarus elaboration, Yosys elaboration, positive-seed Icarus simulations, selected bounded SymbiYosys proofs, structural QoR aggregation, and benchmark aggregation path are reproducible.

\begin{table}[t]
  \centering
  \caption{Seed artifact smoke results}
  \begin{tabular}{lll}
    \toprule
    Metric & Result & Scope \\
    \midrule
    Expected outcome & 57/57 & all seed tasks \\
    Compose pass & 31/31 & positive tasks \\
    Lint/elab pass & 31/31 & positive tasks \\
    Simulation pass & 4/4 & harness-enabled positives \\
    Selected formal pass & 2/2 & direct + width seeds \\
    Structural QoR & 4/4 & reference-enabled positives \\
    Unsafe rejection & 26/26 & negative tasks \\
    JSON AST path & 57/57 & source-to-AST checks \\
    \bottomrule
  \end{tabular}
\end{table}

The paper therefore does not claim final LLM pass-rate improvements yet. The CCF-A submission target still requires running multiple LLM baselines, extending simulation and formal harnesses, adding dedicated subsystem case studies, and reporting richer QoR deltas for adapter-heavy tasks. The value of the current result is reproducibility: a clean checkout can run the Rust tests, emit wrappers and SVA skeletons, reject unsafe negative examples, execute supported positive-seed simulations, prove the selected direct and width seed properties, aggregate machine-readable benchmark JSON including structural QoR CSV/TeX summaries, and pass open-source EDA smoke checks in a controlled Docker environment.

\section{Artifact Reproducibility}

MICO is packaged as a reproducible research artifact rather than a prompt-only demonstration. Rust compilation, unit tests, open-source RTL checks, benchmark tasks, provider smoke tests, and paper builds run through documented repository commands. The reproduced validation for this revision uses the Ubuntu 24.04 Docker EDA environment for Rust, Python, open-source RTL tools, benchmark aggregation, and LLM provider checks. Vivado is not required for the deterministic benchmark results, but a representative host-Vivado subset is available for Xilinx-specific out-of-context synthesis evidence. The paper PDF is built with the Windows-host LaTeX toolchain because the current Docker image does not include LaTeX.

\begin{table}[t]
  \centering
  \caption{Current reproducibility gates}
  \begin{tabular}{lll}
    \toprule
    Gate & Evidence & Status \\
    \midrule
    Rust & fmt/check/test & pass \\
    EDA & Verilator/Icarus/Yosys/SBY & pass \\
    Bench & dev + held-out & pass \\
    LLM & batch validate/offline & pass \\
    Paper & host latexmk/pdflatex & pass \\
    Vivado & four-task OOC subset & pass \\
    \bottomrule
  \end{tabular}
\end{table}

The release candidate checklist is:

\begin{itemize}
  \item freeze the Docker image recipe, Rust toolchain, prompts, provider profiles, and benchmark manifests;
  \item run Rust format, check, and tests from a clean checkout;
  \item run Verilator, Icarus, Yosys, and SymbiYosys smoke checks plus public-development and held-out ModuleComposeBench scoring and aggregate table generation;
  \item run the representative Vivado subset through the pinned Windows-host launcher when reporting Xilinx-specific QoR/timing evidence;
  \item run LLM provider validation and batch-runner validate/offline checks with a redacted local configuration and record model/profile identifiers;
  \item compile the paper with the documented Windows-host LaTeX toolchain, or add LaTeX to the Docker image before claiming a Docker PDF build;
  \item verify that generated PDFs, build outputs, local YAML files, logs, and credentials are absent from the staged source tree.
\end{itemize}

The core reproduction commands are:

\begin{itemize}
  \item Benchmark scoring: \texttt{benchmarks/run\_bench.py}.
  \item Paper table generation: \texttt{benchmarks/aggregate\_results.py}.
  \item LLM validate/offline runs: \texttt{scripts/run\_llm\_bench.py}.
  \item Vivado subset: \path{scripts/vivado-qor-subset.tcl}.
  \item Paper PDF build: \texttt{latexmk} on \texttt{paper/main.tex}.
\end{itemize}

For final submission, this checklist must be extended with full paid LLM baseline results, broader simulation/formal harness outputs, richer all-task QoR reports, archived held-out prompts/results, and archived prompts/responses for every evaluated model profile.

\section{Limitations and Threats to Validity}

MICO requires interface and adapter libraries, and an overly restrictive type system can reduce useful model search space. The current v0 contract checker is a ready/valid subset, not arbitrary LTL. The 62 deterministic tasks and held-out split provide useful coverage but do not by themselves support multi-model statistical claims. Simulation covers all positives, but only 20 use task-specific directed harnesses; the remaining 16 use generated ready/valid smoke harnesses. Formal smoke covers 31 single-clock positives, with 14 directed monitors and 17 generated monitors; CDC proof remains unclaimed.

QoR evidence is limited to structural area/wire counts, generic mapped-cell counts against small references, and a four-task Vivado out-of-context measurement subset. It is not routed timing closure, technology-mapped delay for the full benchmark, or board-level implementation. The authenticated low-cost LLM matrix is negative under current prompts/models and should not be generalized to stronger models or future structured-output settings. Provider results may vary across model versions and API settings, so prompts, model identifiers, hashes, compiler/EDA results, costs, and generated artifacts must be logged.

\section{Conclusion}

MICO reframes LLM-assisted RTL integration as compiler-checked module composition. Its core objects are interfaces, contracts, clock domains, adapters, and compose graphs. By constraining model output and validating every candidate through a Rust compiler and EDA tools, MICO aims to make LLM-assisted hardware integration more diagnosable, repairable, and trustworthy than direct free-form RTL generation.

The present artifact demonstrates the deterministic end-to-end path: parse and check MICO source, build typed IR, emit traceable SystemVerilog wrappers, run 57 ModuleComposeBench seed tasks, run seed Verilator/Yosys smoke checks, execute supported positive-seed simulations, prove selected direct and width seed properties, summarize structural QoR, and validate an SDK-backed LLM provider configuration. The remaining work for a strong submission is clear: run model/baseline comparisons, report repair-turn data, add dedicated subsystem case studies, extend formal and QoR coverage, and package the artifact for independent reproduction.


\bibliographystyle{IEEEtran}
\bibliography{related_work}

\end{document}
